跨越语言的障碍，实现不同语言人们之间的自由交流，是人类自古以来的一个梦想。机器翻译理论的研究目的在于应用计算机作为智能处理工具，实现异种自然语言间的自动翻译过程，其技术意义和社会意义都是十分深远的。

然而由于自然语言的复杂性，直至今天机器翻译的研究仍面临着巨大的困难。除词汇歧义和转换变异映射外，结构歧义一直是机器翻译研究中的主要困难之一。这是因为，各种机器翻译方法，无论是基于知识的还是基于经验的，都或多或少地依赖于源语的结构标注信息来完成语言的转换生成过程。


\cnkeywords{自然语言处理；词汇语义驱动；结构消歧；机器翻译；随机语言模型；知识获取}